% Manuscript insertion candidate: peak-dominated hazard asymptotic.

\section*{Peak-dominated cycle-hazard asymptotic}

The closed reduced survivor law permits an additional high-barrier approximation for sinusoidal loading. Let
\begin{equation}
q(t)=q_m+q_a\sin(2\pi f t),
\qquad q_a>0,
\end{equation}
and let $\lambda_p$ be the quasistatic stable spacing at the tensile peak $q_p=q_m+q_a$. Define
\begin{equation}
B=\frac{EA_ca_0}{k_BT},
\qquad
g_p=\Delta\psi_c(\lambda_p),
\qquad
\nu_p=\frac{\sqrt{\phi''(\lambda_p)}}{2\pi t_0}.
\end{equation}

For the effective-potential climb $g(q)$, stationarity of the stable branch gives the exact derivative
\begin{equation}
\frac{dg}{dq}=-(\lambda_c-\lambda_s),
\end{equation}
and hence
\begin{equation}
\frac{d^2g}{dq^2}=\frac{1}{\phi''(\lambda_s)}>0.
\end{equation}
Near the tensile maximum,
\begin{equation}
q(t)=q_p-\frac12q_a(2\pi f)^2(t-t_p)^2+O[(t-t_p)^4],
\end{equation}
so the rare-event cycle integral is Gaussian to leading order. The resulting one-cycle hazard is
\begin{equation}
\boxed{
\mathcal H_c
\sim
\frac{\nu_p}{f\sqrt{2\pi Bq_a(\lambda_c-\lambda_p)}}
\exp(-Bg_p).
}
\end{equation}
For a narrow structural initial state,
\begin{equation}
S_N\simeq\exp(-N\mathcal H_c),
\end{equation}
so
\begin{equation}
N_p
\sim
\frac{-\ln(1-p)f}{\nu_p}
\sqrt{2\pi Bq_a(\lambda_c-\lambda_p)}
\exp(Bg_p).
\end{equation}
This produces an S--N-like stress sensitivity without introducing a Basquin exponent or a fitted life distribution.

The same expression gives a direct temperature-identifiability transform. Writing
\begin{equation}
C=\frac{EA_ca_0}{k_B},
\end{equation}
we obtain
\begin{equation}
\boxed{
\ln\left(\frac{f\mathcal H_c}{\sqrt T}\right)
\sim
\mathrm{const}
-\frac{Cg_p}{T}.
}
\end{equation}
Thus, if a local first-passage hazard can be measured under a fixed stress waveform,
\begin{equation}
\boxed{
A_c
\simeq
-\frac{k_B}{Ea_0g_p}
\frac{d}{d(1/T)}
\ln\left(\frac{f\mathcal H_c}{\sqrt T}\right).
}
\end{equation}
This is an identifiability relation rather than a present calibration.

For the current diagnostic point $100\pm100$ MPa, 20 Hz, 300 K and sensitivity value $A_c/A_0=50$, direct cycle quadrature gives $\mathcal H_c\approx2.298\times10^{-6}$ and the leading peak approximation gives $2.118\times10^{-6}$, an error of about $-7.8\%$. Across the audited rare-event sweep the leading approximation is within roughly 4--15\% of direct quadrature. A synthetic 260--340 K sweep generated with $A_c/A_0=50$ returns $A_c/A_0\approx50.13$ under the transformed-slope inversion.

These results provide a compact analytical prediction and a future experimental route to $A_c$, but they do not permit direct use of specimen S--N data until local-to-specimen correlation scaling is independently established.
